\documentclass[11pt]{article}

\usepackage[a4paper,margin=1in]{geometry}
\usepackage{amsmath,amssymb,amsthm,mathtools}
\usepackage{bm}
\usepackage{mathrsfs}
\usepackage{hyperref}
\usepackage{enumitem}
\usepackage{booktabs}
\usepackage{microtype}

\hypersetup{
  colorlinks=true,
  linkcolor=blue,
  citecolor=blue,
  urlcolor=blue
}

\newtheorem{theorem}{Theorem}
\newtheorem{proposition}[theorem]{Proposition}
\newtheorem{lemma}[theorem]{Lemma}
\newtheorem{corollary}[theorem]{Corollary}
\theoremstyle{definition}
\newtheorem{definition}[theorem]{Definition}
\newtheorem{assumption}[theorem]{Assumption}
\newtheorem{remark}[theorem]{Remark}
\newtheorem{conjecture}[theorem]{Conjecture}

\newcommand{\R}{\mathbb{R}}
\newcommand{\E}{\mathbb{E}}
\newcommand{\Prob}{\mathbb{P}}
\newcommand{\KL}{\mathrm{KL}}
\newcommand{\Ent}{\mathrm{Ent}}
\newcommand{\dom}{\mathrm{dom}}
\newcommand{\supp}{\mathrm{supp}}
\newcommand{\ip}[2]{\left\langle #1, #2 \right\rangle}
\newcommand{\norm}[1]{\left\lVert #1 \right\rVert}
\newcommand{\mc}[1]{\mathcal{#1}}
\newcommand{\ind}{\mathbf{1}}

\title{Theory Skeleton: Mean-Field Training of Two-Layer MLPs\\with Exponential-Family Natural-Parameter Outputs}
\author{Working draft}
\date{\today}

\begin{document}
\maketitle

\begin{abstract}
This note is a paper-oriented theory skeleton for two-layer MLPs that output
distributional parameters. Its main point is simple: for heteroscedastic
Gaussian regression, the standard mean--variance parameterization
$(\mu,\sigma^2)$ makes the negative log-likelihood nonconvex in the network
output, while the natural-parameter parameterization
$\eta=(\eta_1,\eta_2)=(\mu/\sigma^2,-1/(2\sigma^2))$ restores convexity in the
output variable. The draft now focuses on the practice-faithful model: a
two-layer network with a raw scale head trained by ordinary GD/SGD and mapped
to the negative second natural parameter through a softplus link. In that
setting, the transformed output no longer evolves by ordinary gradient flow;
instead, the mean-field limit induces a mobility-weighted flow and, with shared
feature learning, an operator-valued mobility on function space. We isolate the
finite-horizon regularity and stability questions needed for a deterministic
mean-field PDE, and formulate the first flagship goal as a common-path theorem:
for fixed training horizon $[0,T]$, independently trained wide networks with
the same initialization law should both track the same deterministic law
$\rho_t$.
\end{abstract}

\section{Executive Summary}

The project thesis is:
\begin{quote}
For two-layer networks in the mean-field regime, exponential-family
\emph{natural-parameter outputs} are theoretically preferable to direct
mean--variance outputs, because the induced population loss is convex in the
output variable. In the practice-faithful softplus model, this convexity does
not produce a convex lifted objective; instead it determines the output
geometry of the raw-parameter mean-field flow.
\end{quote}

More concretely:
\begin{enumerate}[label=(\roman*)]
  \item the heteroscedastic Gaussian loss in $(\mu,\sigma^2)$ is jointly
  nonconvex;
  \item the same loss in natural coordinates $\eta$ is convex because it is the
  canonical exponential-family negative log-likelihood;
  \item the practical model should be written with a raw scale field $s$ and a
  softplus link $\eta_2=-\psi_{\rm sp}(s)$, so that the optimized variables are
  the raw particle parameters actually used in practice;
  \item Euclidean GD/SGD in those raw parameters pushes forward to a
  mobility-weighted output flow rather than to ordinary gradient flow in the
  transformed parameter;
  \item in full feature learning that mobility is an operator-valued kernel,
  not a pointwise scalar weight;
  \item the first real theorem is therefore not a long-time convergence claim
  but a finite-horizon mean-field approximation and common-path result.
\end{enumerate}

\paragraph{Flagship first target.}
The right first flagship theorem is not a full $t\to\infty$ selection result.
It is a finite-horizon mean-field approximation theorem: for every fixed
$T<\infty$ and every common initialization law $\rho_0$, independently trained
width-$N$ networks should both remain close, on $[0,T]$, to the same
deterministic trajectory $\rho_t$ as $N\to\infty$. This is the precise bridge
needed for OT-based run-to-run comparisons and mode-connectivity arguments.

\paragraph{Status of claims.}
In this draft, Propositions~\ref{prop:nonconvex-mu-v},
\ref{prop:expfam-convexity}, \ref{prop:softplus-domain},
\ref{prop:gaussian-smoothness}, \ref{prop:function-space-mobility}, and
\ref{prop:operator-mobility} are complete. In addition,
Theorem~\ref{thm:complete-full-batch} and
Corollary~\ref{cor:complete-two-run} now give a complete finite-horizon result
in the narrower continuous-time full-batch regime. The remaining open target is
Theorem~\ref{thm:softplus-finite-horizon} (the broader SGD template) and its
two-run Corollary~\ref{cor:two-run-common-path}. The central next task is to
upgrade the full-batch verification arguments to the discrete-time/SGD setting.
Conjecture~\ref{conj:finite-width} remains the main beyond-first-paper bridge.

\section{Setup}

Let $(X,Y)\sim P$ with $X\in\mc{X}\subseteq \R^d$ and $Y\in\mc{Y}$. Consider a
regular exponential family
\begin{equation}
  p_\eta(y) = h(y)\exp\!\big(\ip{\eta}{T(y)} - A(\eta)\big),
  \qquad \eta\in \mc{H}:=\dom A \subset \R^m,
\end{equation}
where $T:\mc{Y}\to \R^m$ is the sufficient statistic and $A$ is the log-partition
function. The conditional negative log-likelihood is
\begin{equation}
  \ell(\eta;y) = A(\eta) - \ip{\eta}{T(y)} - \log h(y).
\end{equation}
Since the term $-\log h(y)$ does not affect optimization over $\eta$, we will
usually drop it.

\begin{remark}
This draft no longer treats sign-constrained affine heads as a parallel main
route. They remain a useful background comparison, but the actual theorem
targets below are written directly for the raw-particle softplus model trained
by ordinary GD/SGD.
\end{remark}

\subsection{Primary Practical Model: A Softplus Raw-Scale Head}

From this point on, the practice-faithful object of interest is a two-layer
network with a raw scale head. Let $\Xi\subset \R^p$ be a parameter domain and
let $\varphi:\Xi\times \mc{X}\to \R$ be a hidden feature map. For
$\theta=(a,c,\xi)\in \R^2\times \Xi$, define
\[
  \Phi_{\rm mean}(\theta,x):=a\,\varphi(\xi,x),
  \qquad
  \Phi_{\rm scale}(\theta,x):=c\,\varphi(\xi,x).
\]
Fix a floor parameter $\lambda_{\min}>0$ and let
\[
  \psi_{\rm sp}(s):=\lambda_{\min}+\log(1+e^s).
\]
At finite width,
\begin{align}
  \eta_{1,N}(x)
  &:=
  \frac1N\sum_{i=1}^N \Phi_{\rm mean}(\theta_i,x),
  \\
  s_N(x)
  &:=
  \frac1N\sum_{i=1}^N \Phi_{\rm scale}(\theta_i,x),
  \\
  \eta_{2,N}(x)
  &:=
  -\psi_{\rm sp}(s_N(x)).
  \label{eq:softplus-finite-width-head}
\end{align}
The corresponding mean-field outputs are
\begin{align}
  \eta_{1,\rho}(x)
  &:=
  \int_{\R^2\times\Xi}\Phi_{\rm mean}(\theta,x)\,\rho(d\theta),
  \\
  s_\rho(x)
  &:=
  \int_{\R^2\times\Xi}\Phi_{\rm scale}(\theta,x)\,\rho(d\theta),
  \\
  \eta_{2,\rho}(x)
  &:=
  -\psi_{\rm sp}(s_\rho(x)).
  \label{eq:softplus-mf-head}
\end{align}

\begin{remark}
This is the intended practical setting. The optimized variables are the raw
particle parameters $(a_i,c_i,\xi_i)$, exactly as in ordinary training; the
softplus nonlinearity appears only in the forward map that converts the raw
scale field into a negative second natural parameter. The later mobility
equations are therefore a pushed-forward description of raw-parameter GD/SGD,
not a different optimization procedure on transformed parameters.
\end{remark}

\begin{proposition}[Automatic negativity of the softplus scale head]
\label{prop:softplus-domain}
For every probability measure $\rho$ on $\R^2\times\Xi$ and every $x\in\mc{X}$,
\[
  \eta_{2,\rho}(x)\le -\lambda_{\min}<0.
\]
The same pointwise bound holds for every finite-width realization
\eqref{eq:softplus-finite-width-head}.
\end{proposition}

\begin{proof}
By definition, $\psi_{\rm sp}(s)\ge \lambda_{\min}$ for every $s\in\R$.
\end{proof}

\begin{proposition}[Gaussian log-partition smoothness on strips]
\label{prop:gaussian-smoothness}
Fix $M>0$ and $\varepsilon>0$. On the strip
  \[
  D_{M,\varepsilon}
  =
  \{(\eta_1,\eta_2)\in \R^2: |\eta_1|\le M,\; \eta_2\le -\varepsilon\},
  \]
the Gaussian log-partition function
\[
  A_{\rm G}(\eta_1,\eta_2)
  =
  -\frac{\eta_1^2}{4\eta_2}
  + \frac12 \log\!\Big(\frac{-\pi}{\eta_2}\Big)
\]
has globally Lipschitz gradient. More precisely,
\begin{equation}
  \sup_{(\eta_1,\eta_2)\in D_{M,\varepsilon}}
  \norm{\nabla^2 A_{\rm G}(\eta_1,\eta_2)}_{\mathrm{op}}
  \le
  L_{\rm G}(M,\varepsilon),
\end{equation}
where one may take
\begin{equation}
  L_{\rm G}(M,\varepsilon)
  :=
  \frac{1}{2\varepsilon}
  +
  \frac{M}{2\varepsilon^2}
  +
  \frac{M^2}{2\varepsilon^3}
  +
  \frac{1}{2\varepsilon^2}.
\end{equation}
\end{proposition}

\begin{proof}
The Hessian entries are
\begin{align}
  \partial_{11}^2 A_{\rm G}(\eta)
  &= -\frac{1}{2\eta_2},
  \\
  \partial_{12}^2 A_{\rm G}(\eta)
  &= \frac{\eta_1}{2\eta_2^2},
  \\
  \partial_{22}^2 A_{\rm G}(\eta)
  &=
  -\frac{\eta_1^2}{2\eta_2^3}
  +
  \frac{1}{2\eta_2^2}.
\end{align}
On $D_{M,\varepsilon}$ we have
\[
  \big|\partial_{11}^2 A_{\rm G}\big| \le \frac{1}{2\varepsilon},
  \qquad
  \big|\partial_{12}^2 A_{\rm G}\big| \le \frac{M}{2\varepsilon^2},
  \qquad
  \big|\partial_{22}^2 A_{\rm G}\big|
  \le
  \frac{M^2}{2\varepsilon^3} + \frac{1}{2\varepsilon^2}.
\]
The operator norm is bounded by any matrix norm dominating it; for instance,
  the maximum absolute row-sum norm gives the displayed constant.
\end{proof}

\section{Why the Mean--Variance Head Is Theoretically Awkward}

For heteroscedastic Gaussian regression with output $(\mu,v)$ where $v>0$, the
negative log-likelihood is
\begin{equation}
  \ell_{\mu v}(\mu,v;y) = \frac12 \log v + \frac{(y-\mu)^2}{2v}.
\end{equation}

\begin{proposition}[Joint nonconvexity in mean--variance coordinates]
\label{prop:nonconvex-mu-v}
For every $y\in \R$ and every $v>0$, the map $(\mu,v)\mapsto
\ell_{\mu v}(\mu,v;y)$ is not convex on $\R\times (0,\infty)$.
\end{proposition}

\begin{proof}
Its Hessian is
\begin{equation}
  \nabla^2 \ell_{\mu v}(\mu,v;y)
  =
  \begin{pmatrix}
    v^{-1} & (y-\mu)v^{-2} \\
    (y-\mu)v^{-2} & -\frac{1}{2}v^{-2} + (y-\mu)^2 v^{-3}
  \end{pmatrix}.
\end{equation}
The determinant is
\begin{equation}
  \det \nabla^2 \ell_{\mu v}(\mu,v;y)
  =
  -\frac{1}{2v^3} < 0.
\end{equation}
Hence the Hessian is indefinite for all $(\mu,v)$, so the loss is nowhere
jointly convex.
\end{proof}

\begin{remark}
This single calculation already explains why the standard mean-field template
\[
  \rho \mapsto R\!\left(\int \Phi\,d\rho\right) + G(\rho)
\]
does not directly yield a convex lifted objective for a mean--variance head:
the outer functional $R$ is itself nonconvex in the predicted pair
$(\mu(\cdot),v(\cdot))$.
\end{remark}

\section{Natural Parameters Restore Output Convexity}
\label{sec:natural-convexity}

\subsection{General exponential-family statement}

\begin{proposition}[Convexity of the canonical negative log-likelihood]
\label{prop:expfam-convexity}
Let $\{p_\eta:\eta\in \mc{H}\}$ be a regular exponential family. Then for each
fixed $y$, the map
\begin{equation}
  \eta \mapsto \ell(\eta;y) = A(\eta) - \ip{\eta}{T(y)}
\end{equation}
is convex on $\mc{H}$. If the family is minimal, then $A$ is strictly convex,
and therefore $\ell(\cdot;y)$ is strictly convex.
\end{proposition}

\begin{proof}
For exponential families, the log-partition function $A$ is convex on its
natural-parameter domain; in minimal families it is strictly convex
\cite{WainwrightJordan2008}. The linear term $-\ip{\eta}{T(y)}$ preserves
convexity and strict convexity.
\end{proof}

\subsection{Gaussian special case}

For a univariate Gaussian with mean $\mu$ and variance $\sigma^2$, define the
natural parameters
\begin{equation}
  \eta_1 = \mu/\sigma^2,
  \qquad
  \eta_2 = -\frac{1}{2\sigma^2},
  \qquad
  \eta_2 < 0.
\end{equation}
With sufficient statistic $T(y)=(y,y^2)$ and the convention $h(y)=1$
(absorbing $1/\sqrt{2\pi}$ into the log-partition function), the density can
be written
\begin{equation}
  p_\eta(y) = \exp\!\big(\eta_1 y + \eta_2 y^2 - A_{\rm G}(\eta)\big),
  \qquad \eta_2<0,
\end{equation}
where
\begin{equation}
  A_{\rm G}(\eta)
  =
  -\frac{\eta_1^2}{4\eta_2}
  + \frac12 \log\!\Big(\frac{-\pi}{\eta_2}\Big).
\end{equation}
Therefore,
\begin{equation}
  \ell_{\eta}(\eta;y)
  =
  A_{\rm G}(\eta) - \eta_1 y - \eta_2 y^2,
\end{equation}
which is convex on the open half-space $\{\eta_2<0\}$.

\begin{remark}
Output-space convexity is the real conceptual takeaway. In an earlier draft one
could combine it with an affine head to obtain a convex lifted objective over
measures, but that route is no longer the main narrative because it does not
match how softplus scale heads are actually trained.
\end{remark}

\section{Softplus Mean-Field Dynamics}

From this point on, the main deterministic object is the practice-faithful
monotone-link dynamics induced by the raw scale variable. The relevant
mean-field theorem should be formulated in raw parameter space and only then
pushed forward through the softplus output map.

\subsection{Monotone-link reparameterization and operator-valued mobility}

If one could use an affine head to map measures directly to natural parameters,
the output-level convexity would lift to a convex objective over measures.
However, practical implementations
usually predict a raw scalar field $s$ and then enforce positivity of a scale
or precision parameter through a monotone link such as $\exp$ or
$\mathrm{softplus}$. This destroys the affine map $\rho\mapsto \eta_\rho$, so
the convex-lifting argument no longer applies directly. The correct object is a
reparameterized gradient flow.

\begin{proposition}[Function-space reparameterized mobility flow]
\label{prop:function-space-mobility}
Let $I\subset (0,\infty)$ be an open interval, let $\psi:\R\to I$ be a $C^1$
diffeomorphism with $\psi'>0$, and let $\mc{R}$ be a Fr\'echet differentiable
functional on $L^2(P_X;I)$. Define
\[
  \mc{J}(s) := \mc{R}(\psi\circ s),
  \qquad s\in L^2(P_X;\R).
\]
If $s_t$ is a smooth solution of the $L^2(P_X)$ gradient flow
\[
  \partial_t s_t = -\frac{\delta \mc{J}}{\delta s}(s_t),
\]
and $\lambda_t:=\psi(s_t)$, then
\begin{equation}
  \partial_t \lambda_t(x)
  =
  -g(\lambda_t(x))
  \frac{\delta \mc{R}}{\delta \lambda}(\lambda_t)(x),
  \label{eq:function-space-mobility}
\end{equation}
where
\[
  g(\lambda)
  :=
  \big(\psi'(\psi^{-1}(\lambda))\big)^2.
\]
Consequently,
\begin{equation}
  \frac{d}{dt}\mc{R}(\lambda_t)
  =
  -
  \E\Bigg[
    g(\lambda_t(X))
    \left|
      \frac{\delta \mc{R}}{\delta \lambda}(\lambda_t)(X)
    \right|^2
  \Bigg]
  \le 0.
  \label{eq:function-space-dissipation}
\end{equation}
If along the trajectory there exist constants $0<g_-\le g_+<\infty$ such that
$g_- \le g(\lambda_t(x))\le g_+$ for $P_X$-almost every $x$ and all $t$, and if
$\mc{R}$ is $\alpha$-strongly convex on a forward-invariant convex set
$\Lambda\subset L^2(P_X;I)$, then $\lambda_t$ converges to the unique minimizer
$\lambda^\star$ of $\mc{R}$ on $\Lambda$ and
\begin{equation}
  \mc{R}(\lambda_t)-\mc{R}(\lambda^\star)
  \le
  e^{-2\alpha g_- t}
  \big(\mc{R}(\lambda_0)-\mc{R}(\lambda^\star)\big).
  \label{eq:function-space-exp-rate}
\end{equation}
\end{proposition}

\begin{proof}
The chain rule gives
\[
  \frac{\delta \mc{J}}{\delta s}(s)
  =
  \left(
    \frac{\delta \mc{R}}{\delta \lambda}(\psi\circ s)
  \right)\psi'(s).
\]
Therefore
\[
  \partial_t \lambda_t
  =
  \psi'(s_t)\partial_t s_t
  =
  -
  \big(\psi'(s_t)\big)^2
  \frac{\delta \mc{R}}{\delta \lambda}(\lambda_t),
\]
which is \eqref{eq:function-space-mobility}. Pairing with
$\frac{\delta \mc{R}}{\delta \lambda}(\lambda_t)$ in $L^2(P_X)$ yields
\eqref{eq:function-space-dissipation}. If $\mc{R}$ is $\alpha$-strongly convex,
then for every $\lambda\in\Lambda$,
\[
  \left\|
    \frac{\delta \mc{R}}{\delta \lambda}(\lambda)
  \right\|_{L^2(P_X)}^2
  \ge
  2\alpha\big(\mc{R}(\lambda)-\mc{R}(\lambda^\star)\big).
\]
Combining this with \eqref{eq:function-space-dissipation} and the lower bound
$g\ge g_-$ yields
\[
  \frac{d}{dt}\big(\mc{R}(\lambda_t)-\mc{R}(\lambda^\star)\big)
  \le
  -2\alpha g_-
  \big(\mc{R}(\lambda_t)-\mc{R}(\lambda^\star)\big),
\]
and Gr\"onwall's inequality proves \eqref{eq:function-space-exp-rate}.
\end{proof}

\begin{remark}
For Gaussian natural parameters it is convenient to write
$\lambda:=-\eta_2>0$. Since the sign flip $\eta_2\mapsto -\lambda$ is affine,
convexity of the Gaussian negative log-likelihood in $(\eta_1,\eta_2)$ on
$\{\eta_2<0\}$ is equivalent to convexity in $(\eta_1,\lambda)$ on
$\{\lambda>0\}$. The loss of convexity arises only when $\lambda$ is written as
a nonlinear link $\psi(s)$ of a raw network output.
\end{remark}

\begin{proposition}[Operator-valued mobility for a raw-scale mean-field head]
\label{prop:operator-mobility}
Let $\Theta\subset \R^q$ and let $\phi:\Theta\times \mc{X}\to \R$ be $C^1$ in
its first argument. For $m\in \mc{P}(\Theta)$ define
\[
  s_m(x) := \int_\Theta \phi(\theta,x)\,m(d\theta),
  \qquad
  \lambda_m(x) := \psi(s_m(x)),
\]
and the mean-field objective
\[
  F_\psi(m) := \mc{R}(\lambda_m),
\]
where $\mc{R}$ is Fr\'echet differentiable on $L^2(P_X;I)$. Assume $m_t$ is a
sufficiently smooth Wasserstein gradient flow of $F_\psi$,
\begin{equation}
  \partial_t m_t
  =
  \nabla_\theta \cdot
  \left(
    m_t \nabla_\theta \frac{\delta F_\psi}{\delta m}(m_t)
  \right),
  \label{eq:reparam-wgf}
\end{equation}
and let
\[
  \lambda_t:=\lambda_{m_t},
  \qquad
  u_t:=\frac{\delta \mc{R}}{\delta \lambda}(\lambda_t).
\]
Then
\begin{equation}
  \partial_t \lambda_t
  =
  -\mathsf{K}_{m_t,\lambda_t}u_t,
  \label{eq:operator-mobility-flow}
\end{equation}
where $\mathsf{K}_{m,\lambda}:L^2(P_X)\to L^2(P_X)$ is the integral operator
\[
  (\mathsf{K}_{m,\lambda}u)(x)
  :=
  \int_{\mc{X}} K_{m,\lambda}(x,x')u(x')\,P_X(dx')
\]
with kernel
\begin{equation}
  K_{m,\lambda}(x,x')
  :=
  \psi'(s_m(x))\psi'(s_m(x'))
  \int_\Theta
  \ip{\nabla_\theta \phi(\theta,x)}{\nabla_\theta \phi(\theta,x')}
  \,m(d\theta).
  \label{eq:operator-mobility-kernel}
\end{equation}
In particular, $\mathsf{K}_{m,\lambda}$ is self-adjoint and positive
semidefinite, and
\begin{equation}
  \frac{d}{dt}\mc{R}(\lambda_t)
  =
  -
  \ip{u_t}{\mathsf{K}_{m_t,\lambda_t}u_t}_{L^2(P_X)}
  \le 0.
  \label{eq:operator-dissipation}
\end{equation}
If there exists $\kappa>0$ such that
\begin{equation}
  \ip{u}{\mathsf{K}_{m_t,\lambda_t}u}_{L^2(P_X)}
  \ge
  \kappa \norm{u}_{L^2(P_X)}^2
  \qquad
  \forall u \in L^2(P_X),\ \forall t\ge 0,
  \label{eq:operator-coercive}
\end{equation}
and if $\mc{R}$ is $\alpha$-strongly convex in $\lambda$, then
\begin{equation}
  \mc{R}(\lambda_t)-\mc{R}(\lambda^\star)
  \le
  e^{-2\alpha \kappa t}
  \big(\mc{R}(\lambda_0)-\mc{R}(\lambda^\star)\big).
  \label{eq:operator-exp-rate}
\end{equation}
\end{proposition}

\begin{proof}
For $m\in \mc{P}(\Theta)$ and $\theta\in\Theta$, the chain rule gives
\[
  \frac{\delta F_\psi}{\delta m}(m,\theta)
  =
  \E\Big[
    \frac{\delta \mc{R}}{\delta \lambda}(\lambda_m)(X)
    \psi'(s_m(X))
    \phi(\theta,X)
  \Big].
\]
Hence
\[
  \nabla_\theta \frac{\delta F_\psi}{\delta m}(m,\theta)
  =
  \E\Big[
    \frac{\delta \mc{R}}{\delta \lambda}(\lambda_m)(X)
    \psi'(s_m(X))
    \nabla_\theta \phi(\theta,X)
  \Big].
\]
Differentiating $s_{m_t}(x)=\int \phi(\theta,x)\,m_t(d\theta)$ along the
continuity equation \eqref{eq:reparam-wgf} and integrating by parts in $\theta$
yield
\[
  \partial_t s_{m_t}(x)
  =
  -
  \int_\Theta
  \ip{
    \nabla_\theta \phi(\theta,x)
  }{
    \nabla_\theta \frac{\delta F_\psi}{\delta m}(m_t,\theta)
  }
  \,m_t(d\theta).
\]
Multiplying by $\psi'(s_{m_t}(x))$ gives
\[
  \partial_t \lambda_t(x)
  =
  -
  \psi'(s_{m_t}(x))
  \int_\Theta
  \ip{
    \nabla_\theta \phi(\theta,x)
  }{
    \E\big[u_t(X')\psi'(s_{m_t}(X'))\nabla_\theta \phi(\theta,X')\big]
  }
  \,m_t(d\theta),
\]
which is exactly \eqref{eq:operator-mobility-flow}--\eqref{eq:operator-mobility-kernel}.
The kernel is symmetric, and for every $u\in L^2(P_X)$,
\[
  \ip{u}{\mathsf{K}_{m,\lambda}u}_{L^2(P_X)}
  =
  \int_\Theta
  \left\|
    \E\big[
      u(X)\psi'(s_m(X))\nabla_\theta \phi(\theta,X)
    \big]
  \right\|^2
  m(d\theta)
  \ge 0,
\]
so $\mathsf{K}_{m,\lambda}$ is positive semidefinite. Pairing
\eqref{eq:operator-mobility-flow} with $u_t$ gives
\eqref{eq:operator-dissipation}. Under \eqref{eq:operator-coercive} and strong
convexity of $\mc{R}$, the same argument as in
Proposition~\ref{prop:function-space-mobility} yields
\eqref{eq:operator-exp-rate}.
\end{proof}

\begin{remark}
Proposition~\ref{prop:operator-mobility} is written for a single raw scale head
to isolate the geometry. In a full two-output network with shared features and
raw coordinates for both mean and scale, the same calculation produces a
matrix-valued kernel operator with cross-coupling terms between coordinates.
Thus the practice-faithful deterministic theory is not merely a scalar mobility
flow; it is an operator-valued gradient flow in output space. The main theorem
missing from this draft is a verifiable condition implying the coercivity
\eqref{eq:operator-coercive} along the mean-field trajectory.
\end{remark}

\subsection{What coercivity can realistically mean}

The previous proposition isolates the correct operator, but the raw assumption
\eqref{eq:operator-coercive} is too strong on an infinite-dimensional output
space. The right next step is to factor the operator, identify its active
subspace, and formulate coercivity only where it can actually hold.

\begin{proposition}[Factorization and active subspace]
\label{prop:operator-factorization}
For $(m,\lambda)$ as in Proposition~\ref{prop:operator-mobility}, define the
linear operator
\[
  (\mathsf{A}_{m,\lambda}u)(\theta)
  :=
  \E\big[
    u(X)\psi'(s_m(X))\nabla_\theta \phi(\theta,X)
  \big],
  \qquad
  u\in L^2(P_X).
\]
Then
\begin{equation}
  \mathsf{K}_{m,\lambda}
  =
  \mathsf{A}_{m,\lambda}^\ast \mathsf{A}_{m,\lambda},
  \label{eq:operator-factorization}
\end{equation}
and therefore
\begin{equation}
  \ip{u}{\mathsf{K}_{m,\lambda}u}_{L^2(P_X)}
  =
  \norm{\mathsf{A}_{m,\lambda}u}_{L^2(m;\R^q)}^2.
  \label{eq:operator-energy-factorization}
\end{equation}
If
\[
  \mc{V}_{m,\lambda}
  :=
  \overline{\operatorname{Ran}(\mathsf{A}_{m,\lambda}^\ast)}
  =
  \ker(\mathsf{K}_{m,\lambda})^\perp,
\]
then for every $u\in L^2(P_X)$,
\begin{equation}
  \ip{u}{\mathsf{K}_{m,\lambda}u}_{L^2(P_X)}
  =
  \ip{
    P_{\mc{V}_{m,\lambda}}u
  }{
    \mathsf{K}_{m,\lambda}P_{\mc{V}_{m,\lambda}}u
  }_{L^2(P_X)}.
  \label{eq:operator-projected-energy}
\end{equation}
\end{proposition}

\begin{proof}
The identity \eqref{eq:operator-factorization} is exactly the kernel formula
\eqref{eq:operator-mobility-kernel} written in operator notation. The energy
identity \eqref{eq:operator-energy-factorization} follows by expanding
$\norm{\mathsf{A}_{m,\lambda}u}_{L^2(m;\R^q)}^2$, which reproduces the kernel
quadratic form. Since $\mathsf{K}_{m,\lambda}$ is self-adjoint and positive
semidefinite, the orthogonal decomposition
$L^2(P_X)=\ker(\mathsf{K}_{m,\lambda})\oplus \ker(\mathsf{K}_{m,\lambda})^\perp$
implies \eqref{eq:operator-projected-energy}.
\end{proof}

\subsection{Finite-horizon mean-field approximation is the main target}

For the softplus head, the first theorem should be proved in the \emph{raw
parameter space}, not in the pushed-forward output variables. This is exactly
the regime treated in the classical finite-time mean-field literature:
one shows that the empirical measure of the SGD particles stays close, on every
fixed horizon $[0,T]$, to a deterministic McKean--Vlasov PDE. The
operator-valued mobility formulas above are then read off from that PDE after
pushing the dynamics forward through the output map.

\begin{assumption}[Softplus finite-horizon regime]
\label{ass:softplus-finite-horizon}
Fix a horizon $T<\infty$. Assume:
\begin{enumerate}[label=(FH\arabic*)]
  \item the step size has the form $s_k=\varepsilon \xi(k\varepsilon)$ with
  $\xi$ bounded and Lipschitz;
  \item the hidden feature map $\varphi(\xi,x)$ is bounded, continuously
  differentiable in $\xi$, and its first derivative is bounded and Lipschitz in
  $\xi$, uniformly in $x$;
  \item the data distribution has bounded support, or more generally has the
  bounded / sub-Gaussian regularity needed for the finite-time arguments of
  \cite{MeiMisiakiewiczMontanari2019,SirignanoSpiliopoulos2020};
  \item the initialization law $\rho_0$ for the raw particle parameters has
  compact support, or the corresponding bounded-moment substitute used in those
  works;
  \item the scale link is the softplus head
  $\psi_{\rm sp}(s)=\lambda_{\min}+\log(1+e^s)$ with $\lambda_{\min}>0$;
  \item along the deterministic trajectory on $[0,T]$ one has a strip bound
  \[
    |\eta_{1,\rho_t}(x)|\le M_T
    \qquad\text{for all relevant }x,
  \]
  so that Proposition~\ref{prop:gaussian-smoothness} applies with
  $\varepsilon=\lambda_{\min}$.
\end{enumerate}
\end{assumption}

\begin{theorem}[Finite-horizon softplus mean-field template]
\label{thm:softplus-finite-horizon}
Assume Assumption~\ref{ass:softplus-finite-horizon}. Let
$\hat\rho_k^{\,N}=\frac1N\sum_{i=1}^N\delta_{\theta_i^k}$ be the empirical
measure of width-$N$ SGD particles trained with the softplus raw-scale head, and
let $\rho_t$ denote the deterministic mean-field law started from $\rho_0$.
Then the intended first theorem has three parts:
\begin{enumerate}[label=(\roman*)]
  \item the nonlinear PDE for $\rho_t$ is well posed on $[0,T]$;
  \item the map $\rho_0\mapsto \rho_t$ is stable on $[0,T]$;
  \item there exists an error term $\delta_{N,\varepsilon}(T)\to 0$ as
  $N\to\infty$ and $\varepsilon\downarrow 0$ such that
  \[
    \sup_{k\le T/\varepsilon}
    W_1\!\big(\hat\rho_k^{\,N},\rho_{k\varepsilon}\big)
    \le
    \delta_{N,\varepsilon}(T)
  \]
  with high probability.
\end{enumerate}
In bounded-coefficient regimes, one expects the same qualitative form as in
\cite{MeiMisiakiewiczMontanari2019}: a combination of Monte-Carlo error
$N^{-1/2}$ and discretization / stochastic-gradient error of order
$\sqrt{\varepsilon}$ or $\varepsilon$.
\end{theorem}

\begin{remark}
The proof strategy should follow the standard four-dynamics decomposition:
deterministic PDE, nonlinear independent-particle system, interacting particle
system, and discrete SGD iterates. The only genuinely new work is to verify that
the Gaussian softplus head keeps the drift inside a finite-horizon regularity
class on which those comparison arguments close.
\end{remark}

\begin{corollary}[Two-run common-path corollary]
\label{cor:two-run-common-path}
Under the assumptions of Theorem~\ref{thm:softplus-finite-horizon}, let
$\hat\rho_{A,k}^{\,N}$ and $\hat\rho_{B,k}^{\,N}$ be two independent width-$N$
runs started from the same initialization law $\rho_0$. Then, with high
probability,
\[
  \sup_{k\le T/\varepsilon}
  W_1\!\big(\hat\rho_{A,k}^{\,N},\hat\rho_{B,k}^{\,N}\big)
  \le
  2\,\delta_{N,\varepsilon}(T).
\]
Equivalently, both runs are close to the same deterministic reference path
$\rho_t$ on $[0,T]$.
\end{corollary}

\begin{proof}[Proof sketch]
Apply Theorem~\ref{thm:softplus-finite-horizon} to each run and use the triangle
inequality in Wasserstein distance.
\end{proof}

\begin{remark}
Corollary~\ref{cor:two-run-common-path} is the precise bridge needed for
OT-based run-to-run geometry. It is exactly the type of statement used by
Ferbach et al.\ to connect wide-network training to linear mode connectivity:
once two runs are both close to the same law $\rho_t$, one can compare them as
two empirical samples from that common law.
\end{remark}

\subsection{A complete first theorem in the continuous-time full-batch regime}

The broad SGD theorem above is still the right long-term target, but a
\emph{complete} proof can already be closed in a slightly narrower setting:
continuous-time full-batch gradient flow on the population risk. This still
uses the same practice-faithful softplus head, the same raw particle
parameters, and the same common-path conclusion.

\begin{assumption}[Complete full-batch regime]
\label{ass:complete-full-batch}
Fix a horizon $T<\infty$. Assume:
\begin{enumerate}[label=(CF\arabic*)]
  \item training is continuous-time full-batch gradient flow on the population
  risk, with the usual mean-field scaling;
  \item the data distribution satisfies $|Y|\le Y_\ast$ almost surely;
  \item the feature map $\varphi:\R^p\times\mc{X}\to\R$ is bounded, and
  $\nabla_\xi \varphi$ is bounded and Lipschitz in $\xi$, uniformly in $x$;
  \item the initialization law $\rho_0$ on $\R^2\times\R^p$ has compact support;
  \item the scale head is the softplus link
  $\psi_{\rm sp}(s)=\lambda_{\min}+\log(1+e^s)$ with $\lambda_{\min}>0$.
\end{enumerate}
\end{assumption}

Write
\[
  K_\varphi:=\sup_{\xi,x}|\varphi(\xi,x)|,
  \qquad
  K_{\nabla}:=\sup_{\xi,x}\norm{\nabla_\xi\varphi(\xi,x)},
\]
and let $L_\varphi$, $L_{\nabla}$ be Lipschitz constants in $\xi$ for
$\varphi$ and $\nabla_\xi\varphi$, respectively. For Gaussian natural
parameters define
\begin{align}
  q_1(\eta;y)
  &:=
  \partial_{\eta_1}\ell_\eta(\eta;y)
  =
  -\frac{\eta_1}{2\eta_2}-y,
  \\
  q_2(\eta;y)
  &:=
  \partial_{\eta_2}\ell_\eta(\eta;y)
  =
  \frac{\eta_1^2}{4\eta_2^2}
  -
  \frac{1}{2\eta_2}
  -
  y^2.
\end{align}

\begin{proposition}[Raw drift formula for the softplus Gaussian head]
\label{prop:softplus-raw-drift}
Assume Assumption~\ref{ass:complete-full-batch}. For
$\theta=(a,c,\xi)\in\R^2\times\R^p$ and a probability measure $\rho$, define
\[
  \eta_\rho(x)
  :=
  \bigl(\eta_{1,\rho}(x),\eta_{2,\rho}(x)\bigr),
  \qquad
  h_\rho(x,y)
  :=
  \psi_{\rm sp}'(s_\rho(x))\,q_2(\eta_\rho(x);y).
\]
Then the mean-field-scaled full-batch gradient flow of the population risk is
the interacting particle system
\begin{equation}
  \dot\theta_i^N(t)
  =
  b\bigl(\theta_i^N(t),\rho_t^N\bigr),
  \qquad
  \rho_t^N:=\frac1N\sum_{i=1}^N \delta_{\theta_i^N(t)},
  \label{eq:complete-particle-system}
\end{equation}
with drift $b=(b_a,b_c,b_\xi)$ given by
\begin{align}
  b_a(\theta,\rho)
  &:=
  -
  \E\big[q_1(\eta_\rho(X);Y)\,\varphi(\xi,X)\big],
  \label{eq:drift-a}
  \\
  b_c(\theta,\rho)
  &:=
  \E\big[h_\rho(X,Y)\,\varphi(\xi,X)\big],
  \label{eq:drift-c}
  \\
  b_\xi(\theta,\rho)
  &:=
  -
  \E\Big[
    \big(a\,q_1(\eta_\rho(X);Y)-c\,h_\rho(X,Y)\big)\nabla_\xi\varphi(\xi,X)
  \Big].
  \label{eq:drift-xi}
\end{align}
\end{proposition}

\begin{proof}
Let
\[
  \mc{L}_N(\theta_1,\dots,\theta_N)
  :=
  \E\big[\ell_\eta(\eta_{\rho_N}(X);Y)\big],
  \qquad
  \rho_N=\frac1N\sum_{j=1}^N\delta_{\theta_j}.
\]
Since
\[
  \partial_{\theta_i}\eta_{1,\rho_N}(x)
  =
  \frac1N\bigl(\varphi(\xi_i,x),0,a_i\nabla_\xi\varphi(\xi_i,x)\bigr)
\]
and
\[
  \partial_{\theta_i}\eta_{2,\rho_N}(x)
  =
  -\frac1N\psi_{\rm sp}'(s_{\rho_N}(x))
  \bigl(0,\varphi(\xi_i,x),c_i\nabla_\xi\varphi(\xi_i,x)\bigr),
\]
the chain rule gives
\[
  \nabla_{\theta_i}\mc{L}_N
  =
  \frac1N
  \begin{pmatrix}
    \E\big[q_1(\eta_{\rho_N}(X);Y)\varphi(\xi_i,X)\big]
    \\
    -\E\big[h_{\rho_N}(X,Y)\varphi(\xi_i,X)\big]
    \\
    \E\big[(a_i q_1(\eta_{\rho_N}(X);Y)-c_i h_{\rho_N}(X,Y))
    \nabla_\xi\varphi(\xi_i,X)\big]
  \end{pmatrix}.
\]
Multiplying by the mean-field factor $N$ in continuous-time gradient flow
yields \eqref{eq:complete-particle-system}--\eqref{eq:drift-xi}.
\end{proof}

\begin{proposition}[Finite-horizon confinement]
\label{prop:softplus-confinement}
Assume Assumption~\ref{ass:complete-full-batch}. Let
\[
  A_0:=\sup_{(a,c,\xi)\in\supp(\rho_0)} |a|,
  \qquad
  C_0:=\sup_{(a,c,\xi)\in\supp(\rho_0)} |c|,
  \qquad
  R_0:=\sup_{(a,c,\xi)\in\supp(\rho_0)} \norm{\xi}.
\]
Define deterministic envelopes $A_t$, $C_t$, $R_t$ on $[0,T]$ by
\begin{align}
  \dot A_t
  &=
  K_\varphi
  \left(
    \frac{K_\varphi A_t}{2\lambda_{\min}}+Y_\ast
  \right),
  \qquad
  A_{t=0}=A_0,
  \label{eq:A-envelope}
  \\
  \dot C_t
  &=
  K_\varphi
  \left(
    \frac{K_\varphi^2 A_t^2}{4\lambda_{\min}^2}
    +
    \frac{1}{2\lambda_{\min}}
    +
    Y_\ast^2
  \right),
  \qquad
  C_{t=0}=C_0,
  \label{eq:C-envelope}
  \\
  \dot R_t
  &=
  K_{\nabla}
  \left[
    A_t\left(\frac{K_\varphi A_t}{2\lambda_{\min}}+Y_\ast\right)
    +
    C_t\left(
      \frac{K_\varphi^2 A_t^2}{4\lambda_{\min}^2}
      +
      \frac{1}{2\lambda_{\min}}
      +
      Y_\ast^2
    \right)
  \right],
  \qquad
  R_{t=0}=R_0.
  \label{eq:R-envelope}
\end{align}
Then $A_t$, $C_t$, and $R_t$ remain finite on $[0,T]$. Moreover every
characteristic of the mean-field flow and every particle path of
\eqref{eq:complete-particle-system} started from $\supp(\rho_0)$ satisfies
\[
  |a(t)|\le A_t,
  \qquad
  |c(t)|\le C_t,
  \qquad
  \norm{\xi(t)}\le R_t
  \qquad
  \forall t\in[0,T].
\]
In particular, if
\[
  B_T
  :=
  \{(a,c,\xi): |a|\le A_T,\; |c|\le C_T,\; \norm{\xi}\le R_T\},
  \qquad
  M_T:=K_\varphi A_T,
\]
then $\supp(\rho_t)\subseteq B_T$ on $[0,T]$ and
\[
  |\eta_{1,\rho_t}(x)|\le M_T
  \qquad
  \forall x\in\mc{X},\ \forall t\in[0,T].
\]
\end{proposition}

\begin{proof}
Because $\eta_{2,\rho}(x)\le -\lambda_{\min}$ and
$|\eta_{1,\rho}(x)|\le K_\varphi A_t$, the explicit formulas for $q_1$ and
$q_2$ give
\[
  |q_1(\eta_\rho(x);y)|
  \le
  \frac{K_\varphi A_t}{2\lambda_{\min}}+Y_\ast
\]
and
\[
  |q_2(\eta_\rho(x);y)|
  \le
  \frac{K_\varphi^2 A_t^2}{4\lambda_{\min}^2}
  +
  \frac{1}{2\lambda_{\min}}
  +
  Y_\ast^2.
\]
Since $0<\psi_{\rm sp}'< 1$, the drift formulas
\eqref{eq:drift-a}--\eqref{eq:drift-xi} imply
\[
  |\dot a(t)|\le \dot A_t,
  \qquad
  |\dot c(t)|\le \dot C_t,
  \qquad
  \norm{\dot\xi(t)}\le \dot R_t.
\]
Comparison for scalar ODEs yields the claimed bounds. The support statement
follows because each characteristic or particle remains inside $B_T$, and the
bound on $\eta_{1,\rho_t}$ is immediate from
$|\eta_{1,\rho_t}(x)|\le \int |a|\,|\varphi(\xi,x)|\,\rho_t(d\theta)\le K_\varphi A_T$.
\end{proof}

\begin{proposition}[Drift regularity on the finite-horizon confinement set]
\label{prop:softplus-drift-regularity}
Assume Assumption~\ref{ass:complete-full-batch}. There exists a constant
$L_T<\infty$ such that for all $\theta,\bar\theta\in B_T$ and all probability
measures $\rho,\nu$ supported on $B_T$,
\begin{equation}
  \norm{b(\theta,\rho)-b(\bar\theta,\nu)}
  \le
  L_T\big(\norm{\theta-\bar\theta}+W_1(\rho,\nu)\big).
  \label{eq:full-batch-drift-lipschitz}
\end{equation}
Moreover $\sup_{\theta\in B_T,\ \supp(\rho)\subseteq B_T}\norm{b(\theta,\rho)}
\le L_T$.
\end{proposition}

\begin{proof}
On $B_T$, the functions
\[
  \theta\mapsto a\,\varphi(\xi,x),
  \qquad
  \theta\mapsto c\,\varphi(\xi,x)
\]
are uniformly bounded and Lipschitz in $\theta$, with constants depending only
on $A_T,C_T,K_\varphi,L_\varphi$. Hence for every $x$,
\[
  |\eta_{1,\rho}(x)-\eta_{1,\nu}(x)| + |s_\rho(x)-s_\nu(x)|
  \le
  C_T^{(1)} W_1(\rho,\nu)
\]
for a finite constant $C_T^{(1)}$. Since $\psi_{\rm sp}'$ is bounded and
Lipschitz, and since Proposition~\ref{prop:gaussian-smoothness} gives a uniform
Lipschitz constant for $\nabla_\eta\ell_\eta$ on
$D_{M_T,\lambda_{\min}}$, both maps
\[
  (\rho,x,y)\mapsto q_1(\eta_\rho(x);y),
  \qquad
  (\rho,x,y)\mapsto h_\rho(x,y)
\]
are uniformly bounded and Lipschitz in $\rho$ on the confinement set. Inserting
these bounds into \eqref{eq:drift-a}--\eqref{eq:drift-xi}, and using the
boundedness/Lipschitz continuity of $\varphi$ and $\nabla_\xi\varphi$ on $B_T$,
gives \eqref{eq:full-batch-drift-lipschitz}. The uniform bound on $b$ follows
from the same estimates with $\bar\theta=\theta$ and $\rho=\nu$.
\end{proof}

\begin{theorem}[Complete finite-horizon theorem for full-batch GD]
\label{thm:complete-full-batch}
Assume Assumption~\ref{ass:complete-full-batch}. Then:
\begin{enumerate}[label=(\roman*)]
  \item there exists a unique measure-valued path
  $(\rho_t)_{t\in[0,T]}\subset \mathcal P(B_T)$ solving the continuity equation
  \begin{equation}
    \partial_t \rho_t + \nabla_\theta\cdot\big(\rho_t\,b(\theta,\rho_t)\big)=0,
    \qquad
    \rho_{t=0}=\rho_0;
    \label{eq:complete-full-batch-pde}
  \end{equation}
  \item for any two solutions $(\rho_t)$ and $(\nu_t)$ of
  \eqref{eq:complete-full-batch-pde} with initial data $\bar\rho_0$ and $\bar\nu_0$
  supported on $B_T$,
  \begin{equation}
    \sup_{t\in[0,T]} W_1(\rho_t,\nu_t)
    \le
    e^{2L_T T} W_1(\bar\rho_0,\bar\nu_0);
    \label{eq:complete-pde-stability}
  \end{equation}
  \item if $\rho_t^N$ denotes the empirical measure of the interacting particle
  system \eqref{eq:complete-particle-system} started from
  $\rho_0^N=\frac1N\sum_{i=1}^N\delta_{\theta_i^0}$, then
  \begin{equation}
    \sup_{t\in[0,T]}W_1(\rho_t^N,\rho_t)
    \le
    e^{2L_T T}W_1(\rho_0^N,\rho_0).
    \label{eq:complete-particle-tracking}
  \end{equation}
\end{enumerate}
Consequently, if $\theta_i^0$ are i.i.d.\ with law $\rho_0$, then
\[
  \sup_{t\in[0,T]}W_1(\rho_t^N,\rho_t)\to 0
\]
almost surely, hence also in probability, as $N\to\infty$.
\end{theorem}

\begin{proof}
For a continuous path $m=(m_t)_{t\in[0,T]}$ in $\mathcal P(B_T)$, let
$X_t^m(\theta)$ be the unique solution of the ODE
\[
  \frac{d}{dt}X_t^m(\theta)=b(X_t^m(\theta),m_t),
  \qquad
  X_0^m(\theta)=\theta.
\]
Proposition~\ref{prop:softplus-drift-regularity} implies global existence on
$[0,T]$ inside $B_T$. Define the map
\[
  \Gamma(m)_t := (X_t^m)_\#\rho_0.
\]
If $m,n$ are two paths, then for every coupled pair of initial points
$(\theta,\bar\theta)$,
\[
  \frac{d}{dt}
  \norm{X_t^m(\theta)-X_t^n(\bar\theta)}
  \le
  L_T\norm{X_t^m(\theta)-X_t^n(\bar\theta)}
  +
  L_T W_1(m_t,n_t).
\]
Gronwall yields
\[
  W_1(\Gamma(m)_t,\Gamma(n)_t)
  \le
  L_T e^{L_T t}\int_0^t W_1(m_s,n_s)\,ds.
\]
Hence $\Gamma$ is a contraction on sufficiently short time intervals; iterating
the fixed-point argument gives existence and uniqueness of $\rho_t$ on $[0,T]$,
proving (i).

For (ii), couple two solutions $\rho_t,\nu_t$ through characteristics
$X_t,Y_t$ satisfying
\[
  \dot X_t=b(X_t,\rho_t),
  \qquad
  \dot Y_t=b(Y_t,\nu_t).
\]
Then
\[
  \frac{d}{dt}\E\norm{X_t-Y_t}
  \le
  L_T \E\norm{X_t-Y_t} + L_T W_1(\rho_t,\nu_t)
  \le
  2L_T \E\norm{X_t-Y_t}.
\]
Taking the infimum over all couplings of $(X_0,Y_0)$ gives
\eqref{eq:complete-pde-stability}.

For (iii), the empirical measure $\rho_t^N$ solves
\eqref{eq:complete-full-batch-pde} in the weak sense because, for every smooth
test function $\chi$,
\[
  \frac{d}{dt}\int \chi(\theta)\,\rho_t^N(d\theta)
  =
  \frac1N\sum_{i=1}^N \nabla\chi(\theta_i^N(t))\cdot
  b(\theta_i^N(t),\rho_t^N)
  =
  \int \nabla\chi(\theta)\cdot b(\theta,\rho_t^N)\,\rho_t^N(d\theta).
\]
Applying the stability estimate from (ii) with initial measures $\rho_0^N$ and
$\rho_0$ yields \eqref{eq:complete-particle-tracking}. Finally,
$W_1(\rho_0^N,\rho_0)\to 0$ almost surely for i.i.d.\ samples from a compactly
supported law, so the same holds for the full trajectory by
\eqref{eq:complete-particle-tracking}.
\end{proof}

\begin{corollary}[Complete two-run common-path theorem]
\label{cor:complete-two-run}
Assume Assumption~\ref{ass:complete-full-batch}. Let
$\rho_{A,t}^N,\rho_{B,t}^N$ be two independent particle systems of the form
\eqref{eq:complete-particle-system}, with i.i.d.\ initial particles sampled
from the same law $\rho_0$. Then
\begin{equation}
  \sup_{t\in[0,T]}W_1(\rho_{A,t}^N,\rho_{B,t}^N)
  \le
  e^{2L_T T}
  \Big(
    W_1(\rho_{A,0}^N,\rho_0)+W_1(\rho_{B,0}^N,\rho_0)
  \Big),
  \label{eq:complete-two-run}
\end{equation}
and therefore
\[
  \sup_{t\in[0,T]}W_1(\rho_{A,t}^N,\rho_{B,t}^N)\to 0
\]
almost surely, hence in probability, as $N\to\infty$.
\end{corollary}

\begin{proof}
Apply Theorem~\ref{thm:complete-full-batch} to each run and use the triangle
inequality.
\end{proof}

\begin{theorem}[Natural-softplus is the right Gaussian head for this program]
\label{thm:natural-softplus-justification}
Consider the two standard Gaussian output parameterizations:
\begin{enumerate}[label=(\alph*)]
  \item the mean--variance head $(\mu,v)\in \R\times(0,\infty)$;
  \item the natural head $(\eta_1,\eta_2)\in \R\times(-\infty,0)$, implemented
  in practice through the softplus map
  $\eta_2=-\psi_{\rm sp}(s)$ with $\lambda_{\min}>0$.
\end{enumerate}
Then:
\begin{enumerate}[label=(\roman*)]
  \item the Gaussian negative log-likelihood is nowhere jointly convex in
  $(\mu,v)$;
  \item the same loss is convex in natural coordinates
  $(\eta_1,\eta_2)$ on the half-space $\{\eta_2<0\}$;
  \item under Assumption~\ref{ass:complete-full-batch}, the practice-faithful
  natural-softplus head satisfies the complete finite-horizon common-path
  theorem of Theorem~\ref{thm:complete-full-batch} and
  Corollary~\ref{cor:complete-two-run}.
\end{enumerate}
Consequently, among these two standard Gaussian parameterizations, the
natural-softplus head is the one that simultaneously provides:
\begin{enumerate}[label=(\alph*$'$)]
  \item output-level convex geometry for the Gaussian loss;
  \item automatic admissibility of the second coordinate through a practical
  softplus implementation;
  \item a complete finite-horizon common-path theorem in the continuous-time
  full-batch mean-field regime.
\end{enumerate}
\end{theorem}

\begin{proof}
Item (i) is Proposition~\ref{prop:nonconvex-mu-v}. Item (ii) follows from
Proposition~\ref{prop:expfam-convexity} together with the Gaussian natural
parameterization discussed in Section~\ref{sec:natural-convexity}. Item (iii) is exactly
Theorem~\ref{thm:complete-full-batch} and Corollary~\ref{cor:complete-two-run}.
The concluding statement is therefore immediate from the three items above.
\end{proof}

\begin{remark}
The point of Theorem~\ref{thm:natural-softplus-justification} is not that
mean-field theory suddenly becomes impossible for every other Gaussian head. The
point is narrower and more useful: within the standard comparison between
mean--variance outputs and natural outputs, the natural-softplus head is the
only one that simultaneously preserves the exponential-family convex geometry
and fits the complete finite-horizon common-path theorem proved above.
\end{remark}

\subsection{What should be cited, and what must be proved here}

For the broad SGD theorem, the first paper should \emph{not} attempt to reprove
finite-horizon mean-field approximation from scratch. The clean strategy is:
\begin{enumerate}[label=(C\arabic*)]
  \item cite the existing finite-horizon particle-to-PDE machinery for
  two-layer networks and McKean--Vlasov limits
  \cite{MeiMisiakiewiczMontanari2019,SirignanoSpiliopoulos2020};
  \item verify that the raw drift induced by the Gaussian softplus head falls
  into the regularity class required by those black-box results;
  \item derive Theorem~\ref{thm:softplus-finite-horizon} and then
  Corollary~\ref{cor:two-run-common-path} as the direct consequence relevant for
  run-to-run geometry.
\end{enumerate}

\begin{remark}
What is fully borrowed from the literature is the four-dynamics proof skeleton,
the fixed-horizon empirical-measure convergence mechanism, and the fact that
two independent runs started from the same law are both close to the same
deterministic path once each run is close to $\rho_t$. What remains model
specific is the verification step: explicit raw-particle gradients for the
Gaussian softplus head, finite-horizon strip and moment bounds, and the
regularity of the resulting shared-feature two-output drift.
\end{remark}

\begin{remark}
This is also why the first paper should avoid leaning too heavily on the
square-loss-specific $V+U$ notation used in some mean-field references. For the
distributional softplus model, the safest formulation is a generic drift
$b(\theta,\rho)$ together with finite-horizon Lipschitz and moment estimates.
The proof still follows the Mei-style template, but the verification is carried
out directly at the level of the raw drift rather than through a special square
loss decomposition.
\end{remark}

\subsection{Long-time optimization questions are separate}

The operator-valued mobility viewpoint remains important, but mostly for the
\emph{next} stage of the project: long-time optimization, convergence rates, and
selection among many possible global minimizers. Those questions are genuinely
harder than the finite-horizon mean-field bridge and should not be conflated
with it.

\begin{remark}
On finite-support data, the mobility operator reduces to a link-weighted
empirical tangent Gram matrix. For softplus heads, this suggests a clean
long-time heuristic: as long as the raw scale outputs stay away from the flat
left tail, the mobility does not degenerate catastrophically. Turning that
heuristic into a theorem is a plausible second paper, but it is not required for
the first finite-horizon result.
\end{remark}

\section{What Is Actually New in This Project?}

The previous sections now identify a single main route worth pursuing in the
first paper: the softplus raw-scale head. The novelty is therefore sharply
localized: establish a finite-horizon deterministic mean-field theorem for that
practice-faithful model, connect it to finite width, and validate empirically
that independently trained runs with the same initialization law follow a common
deterministic path $\rho_t$ on $[0,T]$. The point is not to rebuild all of
mean-field theory from scratch, but to verify that this distributional output
model satisfies the abstract hypotheses of the existing finite-horizon
McKean--Vlasov machinery. This is the correct bridge to OT-based run-to-run
comparisons such as linear mode connectivity. Questions about $t\to\infty$
selection and mobility coercivity can be left for later work.

\subsection{Softplus is the right practical parameterization}

For Gaussian regression the natural domain is the half-space
\[
  \mc{H}_{\rm G} = \{(\eta_1,\eta_2)\in \R^2 : \eta_2<0\}.
\]
For the present project, the correct mechanism is the softplus raw-scale
parameterization
\[
  \eta_2(x) = -\psi_{\rm sp}(s_\rho(x)),
  \qquad
  \psi_{\rm sp}(s)=\lambda_{\min}+\log(1+e^s),
\]
because it matches actual training code and enforces $\eta_2<0$ automatically.
The main remaining burden is therefore not sign control but finite-horizon
regularity: prove that the raw-particle dynamics keeps the outputs in a strip on
which the Gaussian loss and its derivatives remain controlled.

\begin{remark}
Constrained affine heads, projected dynamics, and barrier methods remain useful
comparison points, but they are no longer the main route of the draft. The
point of this note is precisely to understand the softplus model without
pretending that the transformed parameter is optimized directly.
\end{remark}

\subsection{Vector-valued extension is modest, but should still be written}

Existing mean-field papers are often written for scalar outputs because that is
the simplest presentation. Here the practical model already has a
two-coordinate output $(\eta_1,\eta_2)$ with shared features, so the final
paper still needs to write the vector-valued extension carefully. Replacing
$L^2(P_X;\R)$ by $L^2(P_X;\R^m)$ should be technically routine, but a final
paper still needs to write this out carefully:
\begin{enumerate}[label=(\arabic*)]
  \item Fr\'echet differentiability of $\mc{R}$ on the admissible function class,
  \item derivation of the matrix-valued mobility operator for shared-feature
  multi-output networks,
  \item local Lipschitz continuity and finite-horizon stability of that
  operator along the mean-field trajectory,
  \item existence of the mean-field flow together with the strip control needed
  for the Gaussian softplus head,
  \item quantitative finite-width approximation and the resulting two-run
  coupling corollary for shared initialization law.
\end{enumerate}

\section{Finite-Width Theory and Research Roadmap}

The infinite-width story alone is not enough for a NeurIPS paper. A credible
submission should also say something quantitative about finite width.

\begin{conjecture}[Long-time tracking beyond the first theorem]
\label{conj:finite-width}
Assume the setting of Theorem~\ref{thm:softplus-finite-horizon} and suppose, in
addition, that along the deterministic trajectory $\rho_t$ the particle
velocity field $v_t(\theta)$ and its Jacobian satisfy a bound of the form
\begin{equation}
  \norm{\nabla_\theta v_t(\theta)}
  \le
  C_0 + C_1 \norm{v_t(\theta)}
\end{equation}
uniformly on the relevant time horizon. Then a finite-width network with
polynomially many particles should track the mean-field flow up to polynomial
times, following the strategy of \cite{GlasgowWuBruna2025}.
\end{conjecture}

\begin{remark}
Conjecture~\ref{conj:finite-width} belongs to the second stage of the program.
The first paper only needs fixed-horizon control; long-time tracking matters
once the finite-horizon theorem and the two-run common-path corollary are in
place.
\end{remark}

\section{A Concrete Paper Plan}

\subsection{Theorem stack for a first submission}

\begin{center}
\begin{tabular}{@{}p{0.2\linewidth}p{0.33\linewidth}p{0.37\linewidth}@{}}
\toprule
Claim & Content & Expected difficulty \\
\midrule
Theorem A & Output-level convexity in natural coordinates and nonconvexity of
the mean--variance head & Easy; already done \\
\addlinespace
Theorem B & Complete continuous-time full-batch theorem: confinement,
regularity, PDE well-posedness, particle tracking, and two-run common path &
Done in the present draft \\
\addlinespace
Theorem C & Finite-horizon SGD approximation on $[0,T]$ via an existing
Mei/Sirignano-style theorem & Main remaining theorem \\
\addlinespace
Theorem D & Two-run common path from a shared initialization law & Moderate;
immediate after Theorem C and directly useful for OT/LMC \\
\addlinespace
Theorem E & Long-time mobility / coercivity analysis for optimization and
selection & Second-paper target; not needed for the first finite-horizon bridge \\
\bottomrule
\end{tabular}
\end{center}

\subsection{Minimal experimental package}

The experimental section should directly test the theory-driven claims:
\begin{enumerate}[label=(E\arabic*)]
  \item \textbf{Synthetic heteroscedastic regression.} Compare a standard
  mean--variance head against the softplus raw-scale natural head. Measure
  optimization stability, NLL, calibration, and failure modes.
  \item \textbf{Width scaling.} Train widths
  $N\in\{128,256,512,1024,4096\}$ and compare finite-width trajectories against a
  particle approximation to the mean-field PDE.
  \item \textbf{Two-run geometry.} Train independent runs from the same
  initialization law and measure Wasserstein distance, matching quality, and
  linear interpolation barriers.
  \item \textbf{Softplus-tail monitoring.} Track the distribution of the raw
  scale field $s_N(x)$ and verify empirically that training stays away from the
  flat left tail on the time horizons where the mean-field approximation is
  accurate.
  \item \textbf{Optional link ablation.} Compare softplus against a localized or
  clipped exponential head to isolate what is special about softplus.
\end{enumerate}

\subsection{What would count as a publishable contribution?}

A plausible NeurIPS-scale claim is:
\begin{quote}
We identify natural-parameter output heads as the correct mean-field
parameterization for uncertainty-aware two-layer regression. In the practical
softplus setting, wide networks trained from the same initialization law should
share a common deterministic mean-field path $\rho_t$ on every fixed training
horizon, opening OT-based analyses of run-to-run geometry such as linear mode
connectivity.
\end{quote}

This is sharper and more defensible than claiming a completely new mean-field
theory from scratch, while staying very close to the way practitioners actually
train heteroscedastic models.

\section{Immediate Next Steps}

\begin{enumerate}[label=(S\arabic*)]
  \item Package Theorem~\ref{thm:complete-full-batch} and
  Corollary~\ref{cor:complete-two-run} as the clean fully proved baseline
  result.
  \item Identify the exact black-box finite-horizon theorem to cite from
  \cite{MeiMisiakiewiczMontanari2019,SirignanoSpiliopoulos2020}, and rewrite our
  first theorem as a reduction to verifying its hypotheses.
  \item Write the raw-particle SGD dynamics for the Gaussian softplus head in a
  form that exposes the induced drift $b(\theta,\rho)$.
  \item Prove the finite-horizon strip and moment bounds needed to keep that
  drift in the required regularity class.
  \item Verify the drift assumptions of the chosen black-box theorem:
  Lipschitz dependence on $\theta$, stability with respect to $\rho$, and the
  corresponding growth bounds.
  \item State the two-run corollary explicitly: common initialization law
  $\rho_0$ implies that two independent width-$N$ runs are both close to the
  same deterministic trajectory $\rho_t$ on $[0,T]$.
  \item Translate that corollary into an OT-based run-comparison statement in
  the style needed for linear mode connectivity.
  \item Implement a 1D synthetic heteroscedastic benchmark with a standard
  mean--variance head and the softplus raw-scale natural head.
  \item Numerically compare optimization failure modes, kernel spectra, and
  width scaling.
\end{enumerate}

\section*{Bibliographic Notes}

The mean-field optimization viewpoint for two-layer networks is established in
\cite{ChizatBach2018,MeiMontanariNguyen2018,RotskoffVandenEijnden2018}. The
finite-horizon approximation of SGD by a deterministic PDE, which is the key
bridge for the present project, is developed quantitatively in
\cite{MeiMisiakiewiczMontanari2019,SirignanoSpiliopoulos2020}. The run-to-run
application to linear mode connectivity is represented by \cite{FerbachEtAl2024}.
For long-time finite-width approximation beyond the classical short-horizon
regime, see \cite{GlasgowWuBruna2025}. The convexity facts for exponential
families are standard and can be found in \cite{WainwrightJordan2008}. For
information geometry and natural gradients, see \cite{Amari1998}.

\begin{thebibliography}{99}

\bibitem{Amari1998}
Shun-ichi Amari.
\newblock Natural gradient works efficiently in learning.
\newblock \emph{Neural Computation}, 10(2):251--276, 1998.

\bibitem{ChizatBach2018}
L\'ena\"ic Chizat and Francis Bach.
\newblock On the global convergence of gradient descent for over-parameterized
models using optimal transport.
\newblock In \emph{Advances in Neural Information Processing Systems}, 2018.

\bibitem{GlasgowWuBruna2025}
Margalit Glasgow, Denny Wu, and Joan Bruna.
\newblock Mean-field analysis of polynomial-width two-layer neural network
beyond finite time horizon.
\newblock In \emph{Proceedings of the Thirty Eighth Conference on Learning
Theory}, 2025.

\bibitem{FerbachEtAl2024}
Damien Ferbach, Baptiste Goujaud, Gauthier Gidel, and Aymeric Dieuleveut.
\newblock Proving linear mode connectivity of neural networks via optimal
transport.
\newblock arXiv:2310.19103, 2024.

\bibitem{MeiMisiakiewiczMontanari2019}
Song Mei, Theodor Misiakiewicz, and Andrea Montanari.
\newblock Mean-field theory of two-layers neural networks: Dimension-free
bounds and kernel limit.
\newblock In \emph{Proceedings of the Thirty-Second Conference on Learning
Theory}, pages 2388--2464, 2019.

\bibitem{MeiMontanariNguyen2018}
Song Mei, Andrea Montanari, and Phan-Minh Nguyen.
\newblock A mean field view of the landscape of two-layer neural networks.
\newblock \emph{Proceedings of the National Academy of Sciences},
115(33):E7665--E7671, 2018.

\bibitem{RotskoffVandenEijnden2018}
Grant M. Rotskoff and Eric Vanden-Eijnden.
\newblock Parameters as interacting particles: Long time convergence and
asymptotic error scaling of neural networks.
\newblock In \emph{Advances in Neural Information Processing Systems}, 2018.

\bibitem{SirignanoSpiliopoulos2020}
Justin Sirignano and Konstantinos Spiliopoulos.
\newblock Mean field analysis of neural networks: A law of large numbers.
\newblock \emph{SIAM Journal on Applied Mathematics}, 80(2):725--752, 2020.

\bibitem{WainwrightJordan2008}
Martin J. Wainwright and Michael I. Jordan.
\newblock Graphical models, exponential families, and variational inference.
\newblock \emph{Foundations and Trends in Machine Learning}, 1(1--2):1--305,
2008.

\end{thebibliography}

\end{document}
